\section{The sample mean and the sample proportion}

The previous section constructed a sampling distribution by listing all possible
samples. This works well for very small models, but the number of possible samples
grows quickly. For larger samples, we often study a statistic through a few
important characteristics of its distribution.

The two most important questions are:
\[
\text{Where is the sampling distribution centred?}
\]
and
\[
\text{How much does the statistic vary from sample to sample?}
\]
Expectation describes the centre, while variance and standard deviation describe
the spread.

\subsection*{The sample mean}

Let
\[
X_1,X_2,\ldots,X_n
\]
be independent observations from the same distribution. Suppose
\[
\E[X_i]=\mu
\qquad\text{and}\qquad
\Var(X_i)=\sigma^2
\]
for every \(i\). The sample mean is
\[
\overline X=\frac1n\sum_{i=1}^n X_i.
\]
It is a random variable because it depends on the random sample.

Using linearity of expectation,
\[
\E[\overline X]
=
\E\left[\frac1n\sum_{i=1}^n X_i\right]
=
\frac1n\sum_{i=1}^n\E[X_i].
\]
Since every observation has expectation \(\mu\),
\[
\E[\overline X]
=
\frac1n(n\mu)
=
\mu.
\]

\begin{theorembox}
If the observations all have expectation \(\mu\), then
\[
\E[\overline X]=\mu.
\]
Thus the sampling distribution of the sample mean is centred at the population
mean.
\end{theorembox}

This does not mean that every observed sample mean equals \(\mu\). Different
samples usually produce different values. The statement concerns the centre of
the entire sampling distribution.

Now consider its variance. Because the observations are independent,
\[
\Var(X_1+\cdots+X_n)
=
\Var(X_1)+\cdots+\Var(X_n).
\]
Therefore
\[
\Var(\overline X)
=
\Var\left(\frac1n\sum_{i=1}^nX_i\right)
=
\frac1{n^2}\sum_{i=1}^n\Var(X_i).
\]
Since every observation has variance \(\sigma^2\),
\[
\Var(\overline X)
=
\frac1{n^2}(n\sigma^2)
=
\frac{\sigma^2}{n}.
\]
Hence
\[
\operatorname{SD}(\overline X)
=
\frac{\sigma}{\sqrt n}.
\]

\begin{theorembox}
For independent observations with common variance \(\sigma^2\),
\[
\Var(\overline X)=\frac{\sigma^2}{n}.
\]
As the sample size increases, the sampling distribution of the mean becomes less
spread out.
\end{theorembox}

The individual observations do not become less variable when \(n\) increases.
The population law has not changed. What becomes less variable is their average.
A single unusually large or small observation has less influence when it is
combined with many other observations.

\begin{examplebox}
Suppose individual observations have
\[
\mu=50,
\qquad
\sigma=12.
\]
For samples of size \(n=4\),
\[
\E[\overline X]=50,
\qquad
\operatorname{SD}(\overline X)=\frac{12}{\sqrt4}=6.
\]
For samples of size \(n=36\),
\[
\E[\overline X]=50,
\qquad
\operatorname{SD}(\overline X)=\frac{12}{\sqrt{36}}=2.
\]
Both sampling distributions are centred at \(50\), but means from samples of
size \(36\) vary less from sample to sample.
\end{examplebox}

At this stage we are describing the centre and spread of the sampling
distribution. We are not yet claiming that the distribution is normal. The shape
of the distribution and normal approximations belong to the next chapter.

\subsection*{The sample proportion}

Suppose each observation records success or failure:
\[
X_i=
\begin{cases}
1,&\text{if the }i\text{-th observation is a success},\\
0,&\text{if the }i\text{-th observation is a failure}.
\end{cases}
\]
If
\[
P(X_i=1)=p,
\]
then
\[
\E[X_i]=p
\qquad\text{and}\qquad
\Var(X_i)=p(1-p).
\]

The sample proportion of successes is
\[
\widehat p
=
\frac{X_1+\cdots+X_n}{n}.
\]
This is exactly the sample mean of zero--one observations. Therefore the general
results for the sample mean apply immediately:
\[
\E[\widehat p]=p
\]
and
\[
\Var(\widehat p)=\frac{p(1-p)}{n}.
\]
Consequently,
\[
\operatorname{SD}(\widehat p)
=
\sqrt{\frac{p(1-p)}{n}}.
\]

\begin{summarybox}
The sample mean and the sample proportion illustrate the same principle:
\[
\text{a statistic varies because the sample varies.}
\]
For the sample mean,
\[
\E[\overline X]=\mu,
\qquad
\Var(\overline X)=\frac{\sigma^2}{n}.
\]
For the sample proportion,
\[
\E[\widehat p]=p,
\qquad
\Var(\widehat p)=\frac{p(1-p)}{n}.
\]
In both cases, increasing the sample size reduces sampling variability without
changing the centre of the distribution.
\end{summarybox}
